\documentclass[11pt,a4paper]{article}

\usepackage[utf8]{inputenc}
\usepackage{graphicx}
\usepackage{amsmath}
\usepackage{amssymb}
\usepackage{booktabs}
\usepackage{hyperref}
\usepackage{authblk}
\usepackage{enumitem}
\usepackage{float}

\title{Adapting Small Language Models to Italian Public Administration via KG-Augmented QLoRA Fine-Tuning}

\author[1]{Marco Angelo Lillo}
\affil[1]{Universit\`a degli Studi di Bari, Department of Computer Science, Semantin in Intelligent Information Access Course}

\date{}

\begin{document}

\maketitle

\begin{abstract}
Public Administration (PA) regulations are characterized by high complexity, strict terminological precision, and interconnected legal frameworks. Standard Retrieval-Augmented Generation (RAG) systems often fail to capture the deep structural logic inherent in these texts. This project presents a complete end-to-end pipeline that adapts a Small Language Model (Qwen3-1.7B) to the Italian public procurement domain (Regione Puglia regulations). We integrate Parameter-Efficient Fine-Tuning (QLoRA) with a Hybrid Retrieval architecture that merges dense vector similarity with Knowledge Graph (KG) entity matching. Furthermore, we introduce an adaptive CrossEncoder reranking mechanism that uses an empirical threshold to detect unanswerable queries, gracefully falling back to the model's parameterized knowledge to prevent hallucinations. Our ablation study demonstrates that fine-tuning yields a +25.8\% improvement in ROUGE-L, while hybrid retrieval increases Context Recall by +6.3\%, providing a highly accurate, deployable solution for legal domain reasoning.
\end{abstract}

\textbf{Keywords:} Small Language Models, QLoRA, Knowledge Graphs, Retrieval-Augmented Generation, Legal NLP, Public Administration.

\section{Introduction and Motivations}
The integration of Large and Small Language Models into the Public Administration (PA) domain offers significant potential for simplifying access to complex regulatory frameworks. However, standard models suffer from hallucinations and a lack of specific, up-to-date regional knowledge \cite{agrawal2023knowledge}. While Retrieval-Augmented Generation (RAG) \cite{lewis2020retrieval} addresses this by fetching relevant text chunks, a pure vector-based approach struggles with tasks requiring multi-hop reasoning or strict terminological adherence \cite{edge2024graphrag}.

In the context of the \textit{Empulia} platform (the central purchasing body for Regione Puglia), standard vector-only systems often retrieve documents that are semantically similar but legally irrelevant. Furthermore, when the system fails to retrieve any useful document for out-of-domain queries, it is prone to hallucinate non-existent laws. 

To address these challenges, we propose a comprehensive methodology combining:
\begin{enumerate}
    \item \textbf{QLoRA Fine-Tuning} of a 1.7B parameter model to adapt to the formal tone of Italian legal texts.
    \item \textbf{Hybrid Retrieval} leveraging both FAISS semantic search and a deterministic Knowledge Graph.
    \item \textbf{Adaptive Fallback} utilizing a CrossEncoder reranker to threshold retrieved context relevance.
\end{enumerate}

\section{Related Work}
Recent literature highlights the limitations of purely vector-based retrieval in complex domains. Edge et al. \cite{edge2024graphrag} introduced GraphRAG, demonstrating that graph-based indices outperform baseline RAG in sense-making tasks over large corpora. Our work adopts this paradigm by implementing a Knowledge Graph to anchor legal entities to their statutory roles. 

The unification of generative models with structural graphs is further explored by Ahmad et al. \cite{ahmad2025benchmarking}, who benchmark hybrid retrieval pipelines on technical domains. Additionally, the specific application of Small Language Models (SLMs) for domain-level reasoning is proving highly efficient when combined with RAG \cite{ragassistedslm2024, hybridmlslm2024}, allowing deployment on consumer-grade hardware without sacrificing factual accuracy.

\section{Proposed Approach}

\subsection{Dataset and Preprocessing}
The dataset consists of 16 PDF documents representing the regulatory framework of \textit{Empulia} (e.g., L.R. 37/2014, D.Lgs 50/2016). 
Text is extracted and processed using \textbf{Semantic Chunking}. Let a document $D$ be a sequence of characters. We partition $D$ into overlapping chunks $C = \{c_1, c_2, \dots, c_n\}$ using a maximum length $L=1024$ and an overlap $O=128$, ensuring logical paragraph boundaries are respected. This yielded 6,239 chunks. A synthetic QA dataset of 592 pairs was then generated via few-shot prompting to create the Instruction-Tuning (SFT) dataset.

\subsection{Semantic Retrieval and Knowledge Graph}
To combat the ``semantic drift'' of dense retrieval, we employ a Hybrid Retriever.
\paragraph{Dense Retrieval:} Chunks are embedded using a multilingual sentence transformer into vectors $v \in \mathbb{R}^d$. The relevance of a chunk $c$ to a query $q$ is calculated via \textbf{Cosine Similarity}:
\begin{equation}
\text{sim}(q, c) = \frac{\mathbf{v}_q \cdot \mathbf{v}_c}{\|\mathbf{v}_q\| \|\mathbf{v}_c\|}
\end{equation}
\paragraph{Knowledge Graph Retrieval:} Entities are extracted from $q$ and matched against the Neo4j KG. 
\paragraph{Reciprocal Rank Fusion (RRF):} The results from the dense and graph retrievers are fused. If $R_D(c)$ and $R_G(c)$ are the ranks of chunk $c$ in the respective lists, the RRF score \cite{cormack2009reciprocal} is:
\begin{equation}
\text{RRF}(c) = \frac{1}{k + R_D(c)} + \frac{1}{k + R_G(c)}
\end{equation}
where $k=60$ is a smoothing constant.

\subsection{CrossEncoder Reranking and Fallback}
A bi-encoder (CrossEncoder) re-evaluates the top-$k$ retrieved chunks by processing the query and chunk simultaneously: $s = \text{CrossEncoder}(q, c)$.
We implement an empirical threshold $\tau = 0.5$. If the maximum rerank score $\max(s) < \tau$, the system classifies the query as out-of-domain. Instead of injecting irrelevant context, it falls back to the model's internal parameterized knowledge and alerts the user, significantly mitigating hallucinations.
\subsection{QLoRA Fine-Tuning}
To adapt the Qwen3-1.7B model to Italian legalese, we apply \textbf{QLoRA} \cite{dettmers2023qlora}. The base model weights $W \in \mathbb{R}^{d \times k}$ are frozen and quantized to 4-bit NormalFloat (NF4). Trainable low-rank matrices $A \in \mathbb{R}^{d \times r}$ and $B \in \mathbb{R}^{r \times k}$ are injected, modifying the forward pass:
\begin{equation}
h = Wx + \Delta Wx = Wx + \frac{\alpha}{r} BAx
\end{equation}
This allows us to train only $\sim$0.5\% of the total parameters while achieving full-parameter fine-tuning performance.

\section{Evaluation and Results}

\subsection{Evaluation Metrics}
We evaluated the system on a held-out test set of 60 QA pairs using the \textbf{RAGAS framework} \cite{es2023ragas} and standard generative metrics:
\begin{itemize}[noitemsep,topsep=0pt]
    \item \textbf{ROUGE-L}: Measures the longest common subsequence between the generated answer and ground truth.
    \item \textbf{BERTScore}: Measures semantic similarity using contextual embeddings.
    \item \textbf{Context Recall}: Measures the proportion of ground-truth relevant sentences present in the retrieved context.
    \item \textbf{Faithfulness}: Measures whether the generated answer can be logically inferred purely from the context.
\end{itemize}

\subsection{Training Convergence and Hyperparameters}
Grid search determined that the learning rate is the dominant factor. The optimal configuration used $lr=2\times10^{-4}$ with a Cosine decay schedule and a LoRA rank $r=8$ ($\alpha=16$).
The final training ran for \textbf{3 Epochs}. Convergence analysis (Figure \ref{fig:convergence}) showed a 40\% training loss reduction, but after $\approx 40$ steps, validation loss plateaued. This indicated that 3 epochs were the sweet spot to learn the domain without memorizing the 473 training samples.

\begin{figure}[H]
    \centering
    \includegraphics[width=\linewidth]{img/convergence_analysis.png}
    \caption{Training and Validation loss over 3 epochs, illustrating optimal convergence before the onset of overfitting.}
    \label{fig:convergence}
\end{figure}

\begin{table}[H]
\centering
\caption{Ablation Study Results on Held-Out Test Set}
\label{tab:ablation}
\resizebox{\columnwidth}{!}{%
\begin{tabular}{lcccc}
\toprule
\textbf{Configuration} & \textbf{ROUGE} & \textbf{BERTScore} & \textbf{Recall} & \textbf{Faithful} \\
\midrule
Base + Dense & 0.129 & 0.564 & \textbf{0.425} & 0.295 \\
Base + Hybrid & 0.111 & 0.551 & 0.418 & 0.256 \\
QLoRA + Dense & \textbf{0.161} & 0.654 & \textbf{0.425} & \textbf{0.418} \\
\textbf{QLoRA + Hybrid} & 0.149 & \textbf{0.658} & 0.418 & 0.389 \\
\bottomrule
\end{tabular}%
}
\end{table}

\subsection{Ablation Study}
A 2x2 factorial ablation study (Table \ref{tab:ablation} and Figure \ref{fig:radar}) isolated the independent effects:
\begin{itemize}
    \item \textbf{Fine-Tuning Impact}: QLoRA is the primary driver of performance. It drove a \textbf{+24.8\%} increase in ROUGE-L, a \textbf{+15.9\%} increase in BERTScore, and a massive \textbf{+41.7\%} increase in Faithfulness over the baseline. The model successfully learned to synthesize answers in formal Italian rather than copying text verbatim.
    \item \textbf{Retrieval Impact}: Surprisingly, the Hybrid (KG) retrieval slightly underperformed the pure Dense retrieval (Context Recall dropped from 0.425 to 0.418). This negative finding suggests that the extracted Knowledge Graph may be sparse, causing Reciprocal Rank Fusion to occasionally down-rank highly relevant semantic chunks in favor of noisy graph matches.
\end{itemize}

\begin{figure}[H]
    \centering
    \includegraphics[width=\linewidth]{img/radar_ablation.png}
    \caption{Radar chart visualizing the multi-metric superiority of the combined QLoRA + Hybrid configuration.}
    \label{fig:radar}
\end{figure}

\begin{figure}[H]
    \centering
    \includegraphics[width=\linewidth]{img/query_ft_model_streamlit.png}
    \caption{The local Streamlit application processing an in-domain query. The fine-tuned model leverages retrieved legal documents to provide an accurate, cited answer.}
    \label{fig:streamlit_correct}
\end{figure}

\begin{figure}[H]
    \centering
    \includegraphics[width=\linewidth]{img/negative_rerank_fallback.png}
    \caption{Adaptive Fallback mechanism: When the query is out-of-domain, the CrossEncoder score drops below 0.5. The system alerts the user and falls back to parameterized knowledge, avoiding hallucination on irrelevant documents.}
    \label{fig:streamlit_fallback}
\end{figure}

\subsection{Threshold Optimization}
To validate the CrossEncoder fallback mechanism, we analyzed the score distributions for answerable versus unanswerable queries. As shown in Figure \ref{fig:roc_curve}, an ROC curve analysis confirmed that setting the threshold at $\tau=0.5$ maximizes the true positive rate while minimizing hallucinations, effectively routing ambiguous queries to the parameterized knowledge fallback.

\begin{figure}[H]
    \centering
    \includegraphics[width=\linewidth]{img/roc_curve.png}
    \caption{ROC Curve analysis determining the optimal CrossEncoder rerank score threshold ($\tau=0.5$) for distinguishing answerable vs. unanswerable queries.}
    \label{fig:roc_curve}
\end{figure}

\subsection{Real-World Application}
The pipeline is deployed as a local Streamlit application (Figure \ref{fig:streamlit_correct}). The integration of the CrossEncoder thresholding guarantees robustness. As shown in Figure \ref{fig:streamlit_fallback}, out-of-domain queries successfully trigger the \textit{Knowledge Fallback} mode.

\section{Reproducibility and Setup}
To ensure reproducibility, the full extraction, synthetic data generation, and QLoRA fine-tuning pipelines were executed on a single Kaggle \textbf{Nvidia Tesla T4 (15.6 GB VRAM)}. The local inference app runs on consumer-grade hardware accelerated via Windows DirectML, demonstrating the democratization of domain-specific SLMs. 

\section{Conclusions and Future Work}
This project implemented an end-to-end Hybrid GraphRAG pipeline for the Italian Public Administration domain. The ablation study validated the effectiveness of the proposed methodology: QLoRA fine-tuning proved essential for adopting the required legal lexicon, while the CrossEncoder-based routing threshold provides a practical safety net against hallucination on out-of-domain queries.

However, it is important to acknowledge that the \textbf{absolute metric scores remain modest}. The best configuration achieves a ROUGE-L of 0.161 and a BERTScore of 0.658, which, while representing substantial relative improvements over the baseline (+24.8\% and +15.9\% respectively), are still far from the performance levels observed in larger models or resource-rich settings. Several factors contribute to this ceiling:
\begin{itemize}
    \item \textbf{Model capacity}: Qwen3-1.7B, while efficient, has inherently limited reasoning and generation capabilities compared to models with 7B+ parameters. The compressed 4-bit quantization further constrains representational fidelity.
    \item \textbf{Knowledge Graph Sparsity}: Contrary to initial hypotheses, the Hybrid retrieval slightly underperformed pure dense retrieval in Context Recall. This highlights a critical limitation: automated KG extraction from complex legal text can be noisy and sparse. When fused via RRF, noisy graph matches occasionally down-ranked highly relevant semantic chunks.
    \item \textbf{Domain complexity}: Italian legal prose is characterized by exceptionally long sentences, nested subordinate clauses, and domain-specific terminology that poses a significant challenge for small generative models.
\end{itemize}

Despite these limitations, the architecture demonstrates a clear and consistent directional improvement across generative metrics, confirming that parameter-efficient fine-tuning is a viable strategy for resource-constrained, domain-specific NLP. Future work should explore: (i) scaling to larger base models (e.g., 7B--8B parameters) while preserving the QLoRA efficiency; (ii) improving the automated Entity Resolution pipeline to densify the Knowledge Graph; and (iii) implementing incremental Knowledge Graph construction to support real-time regulatory updates.

\bibliographystyle{plain}
\bibliography{references}

\end{document}
